\section{Заключение}

В работе представлен программный эмулятор CUDA PTX, решающий все семь
поставленных задач. Ключевые возможности реализации включают: перехват
вызовов CUDA Runtime через механизм \texttt{LD\_PRELOAD} без модификации
бинарного файла приложения; однопроходный инкрементный синтаксический анализатор PTX~ISA~7.x~\cite{PTXISA}
с поддержкой 66~мнемоник и автогенерацией регулярных выражений
из \texttt{instructions.json}; таблицу диспетчеризации с O(1)-поиском
и исполнители для всех реализованных инструкций; алгоритмы \texttt{PDOM\_Diverge}
и \texttt{PDOM\_Merge}, обеспечивающие корректную обработку дивергентных потоков
внутри варпа; профилировщик с четырьмя метриками и HTML-отчётом с аннотацией
строк исходного \texttt{.cu}-файла.
Все 7~интеграционных ядер пройдены при максимальном
отклонении от CPU-эталона не более $10^{-3}$; замедление относительно
V100-PCIE составляет 1\,900--2\,300$\times$ для нетривиальных нагрузок~---
неизбежная стоимость программной эмуляции без JIT-компиляции.
Для задач непрерывной интеграции и однократного анализа узких мест замедление в $1\,900$--$2\,300\times$ приемлемо: прогон эмулятора укладывается в секунды при GPU-времени нетривиальных ядер в сотни микросекунд.

Основное текущее ограничение~--- неполнота покрытия PTX~ISA: ряд инструкций не реализован вовсе (тензорные операции \texttt{mma} и \texttt{ldmatrix}, часть векторных и специальных функций), а у реализованных инструкций поддержаны не все варианты модификаторов и режимов адресации, предусмотренных спецификацией.
Эмулятор прошёл проверку на учебных примерах курса «Программирование
на новых архитектурах» и на лабораторных работах студентов, однако
остаётся первой версией инструмента.
Для практического применения на произвольных промышленных кодах он
требует существенного расширения покрытия инструкций, доработки обработки
крайних случаев и повышения производительности; первым шагом станет
покрытие тензорных инструкций \texttt{mma} и \texttt{ldmatrix}.

\newpage
